%%%%%%%%%%%%%%%%%%%%%%%%%%%%%%%%%%%%%%%%%%%%%%%%%%%%%%%%%%%%%%%%%%%
\renewcommand{\thesection}{Anexos}
\titleformat{\section}[block]{\normalfont\Large\bfseries}{}{0.1em}{\thesection}[]
\renewcommand{\thesubsection}{\Alph{subsection}}
\setcounter{equation}{0} % para que empiece en 1
\renewcommand{\theequation}{\Alph{subsection}.\arabic{equation}}
\setcounter{table}{0} % para que empiece en 1
\renewcommand{\thetable}{\Alph{subsection}.\arabic{table}}
\setcounter{figure}{0} % para que empiece en 1
\renewcommand{\thefigure}{\Alph{subsection}.\arabic{figure}}
%%%%%%%%%%%%%%%%%%%%%%%%%%%%%%%%%%%%%%%%%%%%%%%%%%%%%%%%%%%%%%%%%%%

\section{}

\subsection{Equivalencia entre formalismo tetrádico e de Riemann}
\subsubsection{Conexión afín no formalismo de Cartan}
\subsubsection{Derivada totalmente covariante}
\begin{align}
    \mathfrak{D}_\lambda = \partial_\lambda +\Gamma\indices{^\mu_\lambda_\sigma} + \omega\indices{^a_\lambda_s} - (\leftrightarrow)
\end{align}

\subsubsection{Postulados do vielbein}



%%%%%%%%%%%%%%%%%%%%%%%%%%%%%%%%%%%%%%%%%%%%%%%%%%%%%%%%%%%%%%%%%%%
\setcounter{equation}{0} % para que empiece en 1
\renewcommand{\theequation}{\Alph{subsection}.\arabic{equation}}
\setcounter{table}{0} % para que empiece en 1
\renewcommand{\thetable}{\Alph{subsection}.\arabic{table}}
\setcounter{figure}{0} % para que empiece en 1
\renewcommand{\thefigure}{\Alph{subsection}.\arabic{figure}}
%%%%%%%%%%%%%%%%%%%%%%%%%%%%%%%%%%%%%%%%%%%%%%%%%%%%%%%%%%%%%%%%%%%
\subsection{Apéndice 2}




